\section{Intro}

Metrology — the science of measurements. Most of us spend our entire lives rather oblivious to the fact that a whole branch of science is dedicated to the principle of measurement itself. This, however, is one of the most, if not the most, fundamental branches of science and engineering; without it, neither our theories would be proven nor would our collective and free trade across borders be as smooth as they are today.

Think about this for a moment - have you ever stopped and considered that, over the entire globe, we all agree on the length of 1 cm? This is by no means a coincidence, and furthermore, have you ever thought that 1 cm today is the same as it was 10 years ago? This might sound insane, 1 cm is 1 cm, what is this nutcase on about?

But, alas, this is something that has troubled civilizations for thousands of years. We might make a ruler and declare it the master ruler, make it 1 meter long, and define that this ruler is the definition of 1 m. In fact, this is what we did for many years. Before that, the king's foot was used as a unit of length, as the king was chosen by God, and thus his foot was considered the most accurate — crazy as it might sound. The challenge with a physical ruler, however, is that all materials degrade, and that 1 m might turn out to have changed in 100 years' time.

Armed with this knowledge, we as humans, however, face a problem: what can be trusted? Let me be frank — metrology is, in fact, about questioning everything and quantifying trustworthy numbers that describe how much we do not trust something. This is called \textbf{uncertainty}, and the world is riddled with it. We, as engineers, have a sacred duty to ensure that published data and measurements are of the highest quality, and this falls squarely within the domain of metrology.

Thus, the goal of this document is to introduce engineers to the world of metrology — a field that is both fascinating and, at times, very rigorous. Metrology is the bridge between theory and reality, and an essential tool for every engineer.

